\paragraph{有限分辨率可见层的布尔化、真实规范不变量与逆极限异常障碍}
把附录 \ref{subsec:op-algebra-fold-quantum-channel} 的折叠量子信道、\S\ref{subsec:pom-pw-tannaka-krein-reconstruction} 的 profinite 可见对称、以及 \S\ref{subsec:cdim-hardy-finite-mode-projection} 的 Hardy 有限模态投影拼接后，可把本项目中的“有限分辨率量子性”压缩为一条更细的分层：真实 fold-gauge 不变量先裂解为一个可见相干矩阵扇区与一组暗标量账本；在此之后，折叠可见信道才把该相干扇区压到布尔对角。量子资源仍由纤维多重度场精确控制，而任何真正超出有限层 gauge-commutant 的全局非经典性则只能滞留在 inverse-limit 异常残差中。

\begin{theorem}[有限分辨率可见层的布尔化与量子性的残差外置]\label{thm:conclusion-fold-quantum-visible-booleanization}
\leanverified{paper\_conclusion\_fold\_quantum\_visible\_booleanization}
固定分辨率 \(m\ge 1\)，令 \(\mathcal H_m=\ell^2(\Omega_m)\)，并记
\[
\mathcal H_m=\bigoplus_{x\in X_m}\mathcal H_x,
\qquad
\mathcal H_x:=\Pi_x\mathcal H_m,
\qquad
d_m(x):=\dim \mathcal H_x.
\]
对折叠量子信道
\[
\mathcal E_m(A)=\sum_{x\in X_m}\frac{\Tr(\Pi_xA)}{d_m(x)}\,\Pi_x
\qquad (A\in\mathcal B(\mathcal H_m))
\]
定义可见代数
\[
\mathcal Z_m:=\bigoplus_{x\in X_m}\CC\,\Pi_x\ \cong\ C(X_m).
\]
则：
\begin{enumerate}
  \item \(\mathrm{Im}(\mathcal E_m)=\mathcal Z_m\)
  \item \(\mathcal E_m\) 的全部可见尖事件恰为
  \[
  \Pi_S:=\sum_{x\in S}\Pi_x
  \qquad (S\subseteq X_m)
  \]
  \item 因而有限分辨率可见层的尖事件格规范同构于幂集 \(\mathcal P(X_m)\)，特别地为布尔格。
\end{enumerate}
因此，有限分辨率读出中不存在内生的非布尔量子逻辑；所有非交换性只能居留于被 \(\mathcal E_m\) 抹去的残差/扩张层。
\end{theorem}
\begin{proof}
对任意 \(A\in\mathcal B(\mathcal H_m)\)，由定义立即有
\[
\mathcal E_m(A)=\sum_{x\in X_m}\frac{\Tr(\Pi_xA)}{d_m(x)}\,\Pi_x\in \mathcal Z_m,
\]
故 \(\mathrm{Im}(\mathcal E_m)\subseteq \mathcal Z_m\)。
反之若 \(A=\sum_x a_x\Pi_x\in\mathcal Z_m\)，则
\[
\Tr(\Pi_xA)=a_x\,\Tr(\Pi_x)=a_x\,d_m(x),
\]
从而 \(\mathcal E_m(A)=A\)，故 \(\mathcal Z_m\subseteq \mathrm{Im}(\mathcal E_m)\)。
于是 \(\mathrm{Im}(\mathcal E_m)=\mathcal Z_m\)。
第二条正是命题 \ref{prop:op-algebra-fold-quantum-channel-projection-01-rigidity} 的内容；第三条随之立即成立。证毕。
\end{proof}

\begin{proposition}[射影纯态的规范商恰为经典概率单纯形]\label{prop:conclusion-fold-quantum-projective-simplex-quotient}
\leanverified{paper\_conclusion\_fold\_quantum\_projective\_simplex\_quotient}
令
\[
U_m:=\prod_{x\in X_m}U(\mathcal H_x)
\]
按块对角方式作用于 \(\mathbb P(\mathcal H_m)\)。
则存在规范同构
\[
\mathbb P(\mathcal H_m)/U_m\ \cong\ \Delta(X_m),
\qquad
[\psi]\longmapsto p_\psi,
\qquad
p_\psi(x):=\frac{\|\Pi_x\psi\|^2}{\|\psi\|^2}.
\]
换言之，有限分辨率可见纯态模空间并不产生新的量子几何，而严格退化为经典概率单纯形。
\end{proposition}
\begin{proof}
映射 \( [\psi]\mapsto p_\psi \) 对射影缩放与块内酉作用都不变，故先降到商空间。
任取 \(p=(p_x)_{x\in X_m}\in \Delta(X_m)\)，对每个 \(x\) 任选单位向量 \(\xi_x\in \mathcal H_x\)，令
\[
\psi:=\bigoplus_{x\in X_m}\sqrt{p_x}\,\xi_x.
\]
则 \(\|\psi\|=1\) 且 \(p_\psi=p\)，故满射成立。

若 \(p_\psi=p_\phi\)，归一化后可设 \(\|\psi\|=\|\phi\|=1\)。
记 \(\psi_x:=\Pi_x\psi\)、\(\phi_x:=\Pi_x\phi\)，则对每个 \(x\) 有 \(\|\psi_x\|=\|\phi_x\|\)。
若 \(\psi_x=\phi_x=0\) 则任取 \(U_x\in U(\mathcal H_x)\)；若二者非零，则存在 \(U_x\in U(\mathcal H_x)\) 使 \(U_x\psi_x=\phi_x\)。
令 \(U:=\bigoplus_x U_x\in U_m\)，则 \(U\psi=\phi\)，故二者在 \(U_m\)-商中等价。
因此该映射双射。证毕。
\end{proof}

\begin{theorem}[纤维多重度多重集给出折叠量子信道的有限分辨率完全分类]\label{thm:conclusion-fold-quantum-channel-multiplicity-classification}
\leanverified{paper\_conclusion\_fold\_quantum\_channel\_multiplicity\_classification}
设
\[
\mathcal H=\bigoplus_{x\in X}\mathcal H_x,
\qquad
\mathcal H'=\bigoplus_{y\in Y}\mathcal H'_y
\]
是两组有限维正交分解，并分别定义测量--制备型折叠信道
\[
\mathcal E(A)=\sum_{x\in X}\frac{\Tr(P_xA)}{d_x}\,P_x,
\qquad
\mathcal E'(A')=\sum_{y\in Y}\frac{\Tr(P'_yA')}{d'_y}\,P'_y,
\]
其中 \(P_x,P'_y\) 为块投影，\(d_x=\dim\mathcal H_x\)，\(d'_y=\dim\mathcal H'_y\)。
若 \(\dim\mathcal H=\dim\mathcal H'=n\)，则 \(\mathcal E\) 与 \(\mathcal E'\) 酉共轭，当且仅当其纤维多重度多重集
\[
\mathcal D(\mathcal E):=\{d_x:\ x\in X\},
\qquad
\mathcal D(\mathcal E'):=\{d'_y:\ y\in Y\}
\]
相同。
\end{theorem}
\begin{proof}
若 \(\mathcal D(\mathcal E)=\mathcal D(\mathcal E')\)，则可取双射 \(\sigma:X\to Y\) 使 \(d_x=d'_{\sigma(x)}\)。
对每个 \(x\) 选酉同构 \(V_x:\mathcal H_x\to \mathcal H'_{\sigma(x)}\)，令 \(V:=\bigoplus_x V_x\)。
则 \(VP_xV^\ast=P'_{\sigma(x)}\)，故对任意 \(A\in\mathcal B(\mathcal H)\) 有
\[
V\mathcal E(A)V^\ast=\mathcal E'(VAV^\ast).
\]
从而两信道酉共轭。

反之，若二者酉共轭，则其归一化 Choi 态谱相同。
由定理 \ref{thm:op-algebra-fold-quantum-channel-choi-spectrum-mi}，非零特征值恰为
\[
\lambda_d=\frac{1}{nd},
\]
而若某个整数 \(d\) 在多重度多重集中出现 \(N_d\) 次，则 \(\lambda_d\) 的总代数重数为 \(N_d d^2\)。
因此由 Choi 谱可唯一恢复每个 \(d\) 与其出现次数 \(N_d\)，从而唯一恢复 \(\mathcal D(\mathcal E)\)。
故酉共轭必迫使两多重度多重集相同。证毕。
\end{proof}

\begin{corollary}[时间纤维 Fibonacci 立方谱即量子资源的精确配分函数]\label{cor:conclusion-fold-quantum-fiber-spectrum-resource-partition}
\leanverified{paper\_conclusion\_fold\_quantum\_fiber\_spectrum\_resource\_partition}
若对每个 \(x\in X_m\)，纤维重构图满足
\[
\mathcal F_x\cong \prod_{i=1}^{r(x)}\Gamma_{\ell_i(x)}
\]
（结论 \ref{con:xi-terminal-fiber-reconstruction-cartesian}），则
\[
d_m(x)=|\mathcal F_x|=\prod_{i=1}^{r(x)}F_{\ell_i(x)+2}.
\]
于是折叠量子信道 \(\mathcal E_m\) 的资源量完全由该时间纤维立方谱给出：
\begin{enumerate}
  \item 最小 Stinespring 环境维数为
  \[
  \min\dim(\mathcal K_{\mathrm{env}})
  =
  \sum_{x\in X_m}\prod_{i=1}^{r(x)}F_{\ell_i(x)+2}^{\,2};
  \]
  \item 归一化 Choi 态 \(\rho_{RB}=J(\mathcal E_m)/2^m\) 的正特征值对每个 \(x\) 为
  \[
  \lambda_x
  =
  2^{-m}\prod_{i=1}^{r(x)}F_{\ell_i(x)+2}^{-1},
  \]
  其重数恰为
  \[
  \prod_{i=1}^{r(x)}F_{\ell_i(x)+2}^{\,2};
  \]
  \item Choi 互信息满足
  \[
  I(R\!:\!B)_{\rho_{RB}}=H(\pi_m),
  \qquad
  \pi_m(x)=2^{-m}\prod_{i=1}^{r(x)}F_{\ell_i(x)+2}.
  \]
\end{enumerate}
因此有限分辨率量子资源并非额外原语，而是时间纤维立方谱的精确配分函数。
\end{corollary}
\begin{proof}
由结论 \ref{con:xi-terminal-fiber-reconstruction-cartesian}，\(\mathcal F_x\) 为 Fibonacci 立方的 Cartesian 乘积。
而 \(|\Gamma_\ell|=F_{\ell+2}\)，故
\[
d_m(x)=|\mathcal F_x|=\prod_i |\Gamma_{\ell_i(x)}|
=\prod_i F_{\ell_i(x)+2}.
\]
将该表达代入定理 \ref{thm:op-algebra-fold-quantum-channel-minimal-environment-equals-s2} 与定理 \ref{thm:op-algebra-fold-quantum-channel-choi-spectrum-mi} 即得三式。证毕。
\end{proof}

\begin{corollary}[profinite 可见极限上的有限维量子逻辑不可能性]\label{cor:conclusion-profinite-visible-logic-finite-dim-boolean}
\leanverified{paper\_conclusion\_profinite\_visible\_logic\_finite\_dim\_boolean}
令 \(G_\infty^{\mathrm{BC}}\cong \mathrm{Homeo}_{\mathrm{cyl}}(X_\infty)\) 为定理 \ref{thm:pom-bc-profinite-reconstruction-homeo-cyl} 给出的柱语义保持 profinite 群。
则任意有限维连续表示
\[
\pi:G_\infty^{\mathrm{BC}}\to \GL(V)
\]
都经某一有限层 \(G_m^{\mathrm{BC}}\) 因子化；并且一旦在该有限层施加与折叠读出相容的可见条件期望，其对应的尖事件代数必落在某个交换代数 \(\mathcal Z_m\) 中。
因而其尖事件格必为布尔格。
特别地，inverse-limit 可见理论中的任何真正非布尔量子逻辑，都不能被有限维连续表象捕获。
\end{corollary}
\begin{proof}
定理 \ref{thm:pom-repcont-profinite-colimit} 表明任意有限维连续表示都经某个有限商 \(G_m^{\mathrm{BC}}\) 因子化。
在该有限层上，可见读出由折叠条件期望与投影族 \(\{\Pi_x\}_{x\in X_m}\) 给出，而定理 \ref{thm:conclusion-fold-quantum-visible-booleanization} 已证明相应可见代数恰为交换代数
\(
\mathcal Z_m=\bigoplus_{x\in X_m}\CC\,\Pi_x
\)
且其投影格布尔。
故任何有限维连续扇区都只能看到这一布尔骨架，而不能承载真正的非布尔尖事件逻辑。证毕。
\end{proof}

\begin{proposition}[Born 权重不是附加公理，而是 Hardy 投影与折叠条件期望的派生量]\label{prop:conclusion-born-weight-derived-from-hardy-and-fold-expectation}
\leanverified{paper\_conclusion\_born\_weight\_derived\_from\_hardy\_and\_fold\_expectation}
令 \(\Pi_N\) 为命题 \ref{prop:unit-circle-phase-hardy-kernel-rational} 与 \S\ref{subsec:cdim-hardy-finite-mode-projection} 中的有限模态 Hardy 投影。
对任意 \(f\in L^2(\RR,\dd\mu)\) 令
\[
\psi_N:=\Pi_N f\neq 0,
\qquad
\rho_{N,f}:=\frac{\lvert\psi_N\rangle\langle\psi_N\rvert}{\|\psi_N\|^2}.
\]
则有限分辨率可见概率分布由严格恒等式给出
\[
\mu_{N,m}^f(x)
:=
\Tr\!\bigl(\Pi_x\,\mathcal E_m(\rho_{N,f})\bigr)
=
\frac{\|\Pi_x\psi_N\|^2}{\|\psi_N\|^2}.
\]
又由
\[
(\Pi_N f)(t)=\int_{\RR}K_N(t,y)\,f(y)\,\dd\mu(y)
\]
且 \(K_N\) 为显式有理核（命题 \ref{prop:unit-circle-phase-hardy-kernel-rational}），\(\mu_{N,m}^f(x)\) 遂成为输入振幅 \(f\) 的显式有理积分函数。
因此，在该框架内有限分辨率 Born 权重不是附加公理，而是 \(\Pi_N\) 与 \(\mathcal E_m\) 的派生定理。
\end{proposition}
\begin{proof}
由 \(\Pi_x\Pi_y=\delta_{xy}\Pi_x\) 对任意态 \(\rho\) 有
\[
\Tr\!\bigl(\Pi_x\mathcal E_m(\rho)\bigr)
=
\sum_{y\in X_m}\frac{\Tr(\Pi_y\rho)}{d_m(y)}\,\Tr(\Pi_x\Pi_y)
=
\Tr(\Pi_x\rho).
\]
代入 \(\rho=\rho_{N,f}\) 得
\[
\mu_{N,m}^f(x)
=
\Tr(\Pi_x\rho_{N,f})
=
\frac{\langle \psi_N,\Pi_x\psi_N\rangle}{\|\psi_N\|^2}
=
\frac{\|\Pi_x\psi_N\|^2}{\|\psi_N\|^2}.
\]
最后一条核表示正是命题 \ref{prop:unit-circle-phase-hardy-kernel-rational} 的结论。证毕。
\end{proof}

\begin{theorem}[真实 fold-gauge 不变量代数的矩阵扇区裂解]\label{thm:conclusion-fold-gauge-commutant-matrix-sector}
\leanverified{paper\_conclusion\_fold\_gauge\_commutant\_matrix\_sector}
对固定分辨率 \(m\ge 1\)，令
\[
F_x:=\Fold_m^{-1}(x),
\qquad
d_x:=|F_x|,
\qquad
\mathcal H_m=\ell^2(\Omega_m)=\bigoplus_{x\in X_m}\mathcal H_x,
\qquad
\mathcal H_x:=\ell^2(F_x).
\]
定义纤维置换规范群
\[
G_m^{\mathrm{fib}}:=\prod_{x\in X_m}\Sym(F_x),
\]
并记其在 \(\mathcal H_m\) 上的自然置换表示为 \(\rho_m\)。
对每个 \(x\in X_m\)，令
\[
u_x:=d_x^{-1/2}\sum_{\omega\in F_x}\delta_\omega,
\qquad
P_x^+:=\lvert u_x\rangle\langle u_x\rvert,
\qquad
P_x^-:=\Pi_x-P_x^+.
\]
再定义
\[
Q_m:=\sum_{x\in X_m}P_x^+,
\qquad
U_m:=Q_m\mathcal H_m=\Span\{u_x:\ x\in X_m\},
\qquad
W_m:=\bigoplus_{x\in X_m}P_x^-\mathcal H_m.
\]
则存在自然 \(^*\)-代数同构
\[
\rho_m(G_m^{\mathrm{fib}})'
\cong
\mathcal B(U_m)\ \oplus\ \bigoplus_{\substack{x\in X_m\\ d_x\ge 2}}\CC\,P_x^-.
\]
特别地，
\[
Q_m\,\rho_m(G_m^{\mathrm{fib}})'\,Q_m=\mathcal B(U_m)\cong M_{|X_m|}(\CC).
\]
因此真实 fold-gauge 的固定点代数并不交换；它内生地包含一个维数 \(|X_m|\) 的完整量子矩阵扇区。
\end{theorem}
\begin{proof}
对每个 \(x\in X_m\)，有正交分解
\[
\mathcal H_x=\CC u_x\oplus W_x,
\qquad
W_x:=u_x^\perp=P_x^-\mathcal H_m.
\]
其中 \(\Sym(F_x)\) 在 \(\CC u_x\) 上平凡作用，在 \(W_x\) 上给出标准表示；当 \(d_x\ge 2\) 时，该标准表示在 \(\CC\) 上不可约。
于是
\[
\mathcal H_m
=
U_m\oplus \bigoplus_{\substack{x\in X_m\\ d_x\ge 2}}W_x
\]
是 \(G_m^{\mathrm{fib}}\) 的等型分解，并且 \(U_m\) 是平凡表示的 \(|X_m|\) 重等型块。

由于 \(G_m^{\mathrm{fib}}\) 在 \(U_m\) 上平凡作用，故 \(\mathcal B(U_m)\subset \rho_m(G_m^{\mathrm{fib}})'\)。
另一方面，若 \(x\neq y\) 且 \(T:W_x\to W_y\) 与 \(G_m^{\mathrm{fib}}\) 对易，则把 \(T\) 限制到第 \(x\) 个因子 \(\Sym(F_x)\) 可见：\(W_x\) 上该因子作用非平凡，而 \(W_y\) 上该因子作用平凡，故
\[
T\in \Hom_{\Sym(F_x)}(W_x,\mathbf 1)=0.
\]
同理，
\[
\Hom_{G_m^{\mathrm{fib}}}(U_m,W_x)
=
\Hom_{G_m^{\mathrm{fib}}}(W_x,U_m)
=0.
\]
又由 Schur 引理，
\[
\End_{G_m^{\mathrm{fib}}}(W_x)=\CC\,I_{W_x}=\CC\,P_x^-.
\]
综上，交换子代数恰为
\[
\mathcal B(U_m)\oplus \bigoplus_{\substack{x\in X_m\\ d_x\ge 2}}\CC\,P_x^-.
\]
再乘以 \(Q_m\) 即得
\(
Q_m\rho_m(G_m^{\mathrm{fib}})'Q_m=\mathcal B(U_m)
\)。
证毕。
\end{proof}

\begin{corollary}[真实规范纯态空间复现为射影空间]\label{cor:conclusion-fold-gauge-purestate-projective-space}
\leanverified{paper\_conclusion\_fold\_gauge\_purestate\_projective\_space}
存在规范酉同构
\[
U_m\cong \ell^2(X_m),
\qquad
u_x\longmapsto e_x.
\]
因此真实 fold-gauge 不变量代数中的相干纯态空间同构于
\[
\mathbb P(U_m)\cong \mathbb P^{|X_m|-1}(\CC).
\]
并且
\[
\dim U_m=|X_m|,
\qquad
\dim W_m=\sum_{x\in X_m}(d_x-1)=2^m-|X_m|.
\]
当 \(m=6\) 时，由推论 \ref{cor:fold-bin-m6-fiber-hist} 的直方图
\[
2{:}8,\qquad 3{:}4,\qquad 4{:}9
\]
可得
\[
|X_6|=21,\qquad \dim U_6=21,\qquad \dim W_6=43,
\]
并且
\[
\rho_6(G_6^{\mathrm{fib}})'
\cong
M_{21}(\CC)\oplus \CC^{21}.
\]
\end{corollary}
\begin{proof}
\(\{u_x\}_{x\in X_m}\) 两两正交且均为单位向量，故构成 \(U_m\) 的一组标准正交基，从而给出所示酉同构。
维数公式由
\[
\sum_{x\in X_m}d_x=|\Omega_m|=2^m
\]
立得。
纯态空间结论是有限维 Hilbert 空间的射影化。
最后，当 \(m=6\) 时，由推论 \ref{cor:fold-bin-m6-fiber-hist} 得
\[
|X_6|=8+4+9=21,
\qquad
\dim W_6=8(2-1)+4(3-1)+9(4-1)=43.
\]
再代入定理 \ref{thm:conclusion-fold-gauge-commutant-matrix-sector} 即得
\[
\rho_6(G_6^{\mathrm{fib}})'\cong M_{21}(\CC)\oplus \CC^{21}.
\]
这也解释了命题 \ref{prop:conclusion-fold-quantum-projective-simplex-quotient} 的单纯形并非原初对象，而是对更大的块内酉群 \(\prod_{x\in X_m}U(\mathcal H_x)\) 进一步商化后的结果。证毕。
\end{proof}

\begin{theorem}[可见布尔读出是矩阵扇区退相干与暗多重度坍缩的复合]\label{thm:conclusion-fold-gauge-average-visible-decoherence}
在定理 \ref{thm:conclusion-fold-gauge-commutant-matrix-sector} 的设定下，定义群平均
\[
\mathfrak T_m(A):=\frac{1}{|G_m^{\mathrm{fib}}|}\sum_{g\in G_m^{\mathrm{fib}}}\rho_m(g)A\rho_m(g)^*
\qquad
(A\in \mathcal B(\mathcal H_m)).
\]
则对任意 \(A\in\mathcal B(\mathcal H_m)\)，有显式公式
\[
\mathfrak T_m(A)
=
Q_mAQ_m
+\sum_{\substack{x\in X_m\\ d_x\ge 2}}
\frac{\Tr(P_x^-A)}{d_x-1}\,P_x^-.
\]
另一方面，折叠可见信道满足
\[
\mathcal E_m(A)
=
\sum_{x\in X_m}
\frac{\langle u_x,Au_x\rangle+\Tr(P_x^-A)}{d_x}\,\Pi_x.
\]
因此 \(\mathcal E_m\) 的布尔化并非原初结构，而是先保留整个 \(\mathcal B(U_m)\) 量子矩阵块，再沿基 \(\{u_x\}\) 完全退相干，并把每条纤维的暗标准分量仅以总迹的形式坍缩到相应对角项之后的复合结果。
\end{theorem}
\begin{proof}
群平均 \(\mathfrak T_m\) 是到交换子代数 \(\rho_m(G_m^{\mathrm{fib}})'\) 的条件期望。
由定理 \ref{thm:conclusion-fold-gauge-commutant-matrix-sector}，
\(\rho_m(G_m^{\mathrm{fib}})'\) 在分解
\[
\mathcal H_m=U_m\oplus \bigoplus_{x:d_x\ge 2}W_x
\]
上恰为
\(
\mathcal B(U_m)\oplus \bigoplus_x \CC P_x^-
\)。
故 \(\mathfrak T_m\) 在平凡表示块 \(U_m\) 上不做压缩，保留为 \(Q_mAQ_m\)；
而在每个不可约块 \(W_x\) 上只能留下标量部分，其系数由迹归一化给出：
\[
\frac{\Tr(P_x^-A)}{\dim W_x}
=
\frac{\Tr(P_x^-A)}{d_x-1}.
\]
于是得到第一式。

第二式由
\[
\Pi_x=P_x^++P_x^-,
\qquad
\Tr(P_x^+A)=\langle u_x,Au_x\rangle
\]
直接推出：
\[
\Tr(\Pi_xA)=\langle u_x,Au_x\rangle+\Tr(P_x^-A).
\]
代回 \(\mathcal E_m\) 的定义即得。
因此 \(\mathcal E_m\) 先读取 \(U_m\) 上完整的相干块，再把它退相干到 \(\{u_x\}\) 的对角，并把暗块只保留其总迹权重。证毕。
\leanverified{paper\_conclusion\_fold\_gauge\_average\_visible\_decoherence}
\end{proof}

\begin{theorem}[可见布尔化核的显式分解与精确维数]\label{thm:conclusion-fold-quantum-visible-kernel-explicit-decomposition}
\leanverified{paper\_conclusion\_fold\_quantum\_visible\_kernel\_explicit\_decomposition}
在定理 \ref{thm:conclusion-fold-gauge-average-visible-decoherence} 的设定下，记
\[
A_m^{\mathrm{vis}}:=\rho_m(G_m^{\mathrm{fib}})',
\qquad
E_m^{\mathrm{vis}}:=\mathcal E_m|_{A_m^{\mathrm{vis}}}:A_m^{\mathrm{vis}}\to Z_m,
\]
\[
N_{\ge 2}(m):=\#\{x\in X_m:\ d_x\ge 2\},
\qquad
\mathcal O_m:=\{B\in \mathcal B(U_m):\ \langle u_x,Bu_x\rangle=0\ \forall x\in X_m\}.
\]
则 \(E_m^{\mathrm{vis}}\) 为满射，且其核具有显式直和分解
\[
\ker E_m^{\mathrm{vis}}
=
\mathcal O_m
\oplus
\bigoplus_{\substack{x\in X_m\\ d_x\ge 2}}
\CC\Bigl(P_x^--(d_x-1)P_x^+\Bigr).
\]
特别地，
\[
\dim_{\CC}\ker E_m^{\mathrm{vis}}
=
|X_m|(|X_m|-1)+N_{\ge 2}(m).
\]
因此，可见布尔化在固定分辨率上擦除的是一个精确可计数的非交换核，而不是一团不可分辨的“量子剩余”。
\end{theorem}
\begin{proof}
由定理 \ref{thm:conclusion-fold-gauge-commutant-matrix-sector}，
任意 \(A\in A_m^{\mathrm{vis}}\) 都可唯一写成
\[
A=B+\sum_{\substack{x\in X_m\\ d_x\ge 2}}\lambda_xP_x^-,
\qquad
B\in\mathcal B(U_m).
\]
再由定理 \ref{thm:conclusion-fold-gauge-average-visible-decoherence}，
\[
E_m^{\mathrm{vis}}(A)
=
\sum_{x\in X_m}
\frac{\langle u_x,Bu_x\rangle+(d_x-1)\lambda_x}{d_x}\,\Pi_x,
\]
其中当 \(d_x=1\) 时约定 \(\lambda_x=0\)。

先证满射：对任意
\[
Z=\sum_{x\in X_m}c_x\Pi_x\in Z_m
\]
取
\[
B:=\sum_{x\in X_m}d_xc_x\,P_x^+\in \mathcal B(U_m),
\qquad
\lambda_x:=0,
\]
则 \(E_m^{\mathrm{vis}}(B)=Z\)。

为求核，把 \(B\) 在基 \(\{u_x\}_{x\in X_m}\) 下分解为
\[
B=B_0+\sum_{x\in X_m}b_xP_x^+,
\qquad
B_0\in\mathcal O_m.
\]
于是 \(A\in\ker E_m^{\mathrm{vis}}\) 当且仅当对每个 \(x\in X_m\) 有
\[
b_x+(d_x-1)\lambda_x=0.
\]
故
\[
A
=
B_0+\sum_{\substack{x\in X_m\\ d_x\ge 2}}\lambda_x\Bigl(P_x^--(d_x-1)P_x^+\Bigr),
\]
从而得到所示直和分解。最后，
\[
\dim_{\CC}\mathcal O_m
=
\dim_{\CC}\mathcal B(U_m)-|X_m|
=
|X_m|^2-|X_m|,
\]
每个 \(d_x\ge 2\) 再贡献一个独立一维核方向，故
\[
\dim_{\CC}\ker E_m^{\mathrm{vis}}
=
|X_m|^2-|X_m|+N_{\ge 2}(m).
\]
证毕。
\end{proof}


\leanverified{paper\_conclusion\_fold\_capacity\_wedderburn\_exact\_bit\_threshold}
\begin{corollary}[容量曲线、Wedderburn 型与精确反演门槛的同一性]\label{cor:conclusion-fold-capacity-wedderburn-exact-bit-threshold}
固定分辨率 \(m\ge 1\)。
连续容量曲线
\[
\mathcal C_m^{\mathrm{cont}}(T):=\sum_{x\in X_m}\min(d_x,T)
\qquad (T\ge 0)
\]
唯一决定纤维直方图
\[
n_k(m):=\#\{x\in X_m:\ d_x=k\},
\]
从而唯一决定群胚代数的 Wedderburn 型
\[
\CC[\mathcal R_m^{\mathrm{bin}}]
\cong
\bigoplus_{k\ge 1}M_k(\CC)^{\oplus n_k(m)}.
\]
此外，达到精确反演
\[
\sum_{x\in X_m}\min(d_x,2^B)=2^m
\]
所需的最小侧信息比特数满足
\[
B_{\mathrm{exact}}(m)
=
\Bigl\lceil \log_2 d_{\max}(m)\Bigr\rceil
=
\Bigl\lceil \log_2 \operatorname{Ind}_{\mathrm{PP}}(E_m^Z)\Bigr\rceil,
\]
其中 \(E_m^Z\) 为定理 \ref{thm:xi-time-part9ze-center-condexp-pimsner-popa-index} 的中心条件期望。
特别地，在 window-\(6\) 上
\[
d_{\max}(6)=4,
\qquad
B_{\mathrm{exact}}(6)=2.
\]
\end{corollary}
\begin{proof}
由定理 \ref{thm:xi-time-part60aaa-wedderburn-information-completeness}，
\(\mathcal C_m^{\mathrm{cont}}\)、纤维直方图 \(\{n_k(m)\}_{k\ge 1}\) 与半单代数同构型彼此唯一决定，故第一部分成立。

另一方面，
\[
\sum_{x\in X_m}\min(d_x,2^B)=2^m=\sum_{x\in X_m}d_x
\]
成立，当且仅当 \(2^B\ge d_x\) 对全部 \(x\) 同时成立；
故最小 \(B\) 正是 \(\lceil\log_2 d_{\max}(m)\rceil\)。
再由定理 \ref{thm:xi-time-part9ze-center-condexp-pimsner-popa-index}，
\[
\operatorname{Ind}_{\mathrm{PP}}(E_m^Z)=d_{\max}(m),
\]
遂得第二式。
对 \(m=6\)，由推论 \ref{cor:fold-bin-m6-fiber-hist} 得 \(d_{\max}(6)=4\)，从而 \(B_{\mathrm{exact}}(6)=2\)。证毕。
\end{proof}

\begin{corollary}[bin-fold 的全部正阶 Rényi 熵率坍缩到 \texorpdfstring{$\log\varphi$}{log phi}]\label{cor:conclusion-binfold-renyi-rate-collapse}
\leanverified{paper\_conclusion\_binfold\_renyi\_rate\_collapse}
在二进制区间口径下，令
\[
\mu_m(x):=\frac{d_m^{\mathrm{bin}}(x)}{2^m}
\qquad(x\in X_m),
\]
并对任意 \(q>0\) 定义 Rényi 熵
\[
H_q(\mu_m):=
\begin{cases}
\dfrac{1}{1-q}\log\sum_x \mu_m(x)^q,& q\neq 1,\\[1ex]
-\sum_x \mu_m(x)\log\mu_m(x),& q=1.
\end{cases}
\]
则对一切 \(q>0\) 都有
\[
\lim_{m\to\infty}\frac{1}{m}H_q(\mu_m)=\log\varphi.
\]
因此 bin-fold 的全部正阶 Rényi 熵率共享同一指数斜率，非均匀性只能停留在 \(O(1)\) 的亚指数层。
\end{corollary}
\begin{proof}
当 \(q\neq 1\) 时，定理 \ref{thm:xi-fold-renyi-constant-spectrum-closed-form} 给出
\[
H_q(\mu_m)=m\log\varphi+c(q)+O\!\Bigl(\Bigl(\frac{\varphi}{2}\Bigr)^m\Bigr),
\]
故
\[
\frac{1}{m}H_q(\mu_m)\to \log\varphi.
\]
当 \(q=1\) 时，同一定理给出
\[
H(\mu_m)=m\log\varphi+\frac{\log\varphi}{\varphi\sqrt 5}+o(1),
\]
故极限仍为 \(\log\varphi\)。
于是对全部 \(q>0\) 结论统一成立。证毕。
\end{proof}

\begin{conjecture}[inverse-limit 异常 \(2\)-类是全局量子性的唯一障碍]\label{conj:conclusion-inverse-limit-anomaly-global-quantumity}
记 \(\mathfrak a_\infty^{\mathbf{PW}}\) 为投影词 \(2\)-范畴 \(\mathbf{PW}\) 在 inverse-limit 可见塔上的兼容异常类极限，即由各有限切片中的异常上同调类 \([\mathrm{Anom}]\) 所拼接的残差 \(2\)-类。
则 inverse-limit 可见理论存在真正非布尔的完备尖事件逻辑，当且仅当 \(\mathfrak a_\infty^{\mathbf{PW}}\) 在所有有限截断上都不能被消去。
换言之，若全局层仍保留不可约量子性，则其唯一持续来源应是贯穿全部有限截断而不消失的异常 \(2\)-类，而非任何有限维可见读出本身。
\end{conjecture}

\begin{remark}[支持性理由]\label{rem:conclusion-inverse-limit-anomaly-global-quantumity-support}
定理 \ref{thm:conclusion-fold-quantum-visible-booleanization} 已表明每个有限分辨率可见层都严格布尔化；
推论 \ref{cor:conclusion-profinite-visible-logic-finite-dim-boolean} 进一步说明每个有限维连续对称扇区都只能看到该布尔骨架。
另一方面，\(\mathbf{PW}\) 中的异常是否为零由上同调类控制（推论 \ref{cor:pom-anom-cohomology-class}），而在有界切片上 strictification 的存在性可被归约为有限异常消失约束的可满足性（推论 \ref{cor:pom-strictification-existence-decidable-bounded-slice}）。
特别地，在 window-\(6\) 的最小非平凡锚点上，异常标签可压缩为 \((\ZZ_2)^{21}\) 上的有限线性可满足性问题，并在 boundary 中进一步投影到规范的 \((\ZZ_2)^3\) 子格（推论 \ref{cor:fold-bin-gauge-m6} 与推论 \ref{cor:fold-bin-boundary-center-z2}）。
因此，若 inverse limit 处仍有真正不可压平的非经典性，它就不可能来自任何单个有限层或任何有限维连续扇区，而只能来自这一异常残差类在极限中的持续存活。
\end{remark}

% --- Distillation writeback ---
% --- Distillation writeback: Wang-Zahl fold-aware coherent grain audit ---
\begin{theorem}[fold-aware coherent grains inside the visible matrix sector]\label{distill:wang_zahl-grain_prism_structure_dichotomy-conclusion-visible-matrix-grains}
\textbf{日期时间：2026-04-23T06:15:38+08:00. 依赖状态：rigidity.}
沿用定理 \ref{thm:conclusion-fold-gauge-commutant-matrix-sector} 的记号。对 \(x,y\in X_m\)，令
\[
E_{xy}:=|u_x\rangle\langle u_y|\in\mathcal B(U_m),
\qquad E_{xx}=P_x^+.
\]
设 \(\mathcal A\subset\mathcal B(U_m)\) 是一个 unital \(^\ast\)-子代数，并假设
\[
P_x^+\in\mathcal A\qquad(x\in X_m).
\]
定义图 \(\Gamma_{\mathcal A}\) 的顶点集为 \(X_m\)，且当 \(x\ne y\) 且
\[
P_x^+\mathcal A P_y^+\ne 0
\]
时连接 \(x\) 与 \(y\)。令 \(\mathcal C(\mathcal A)\) 为 \(\Gamma_{\mathcal A}\) 的连通分支，并记
\[
U_C:=\Span\{u_x:x\in C\}.
\]
则
\[
\mathcal A=\bigoplus_{C\in\mathcal C(\mathcal A)}\mathcal B(U_C)
\]
作为 \(U_m=\bigoplus_C U_C\) 上的块对角 \(^\ast\)-代数。

特别地，\(\mathcal C(\mathcal A)\) 正是该 visible audit algebra 的 maximal coherent grains。若 \(Y\in\mathcal C(\mathcal A)\)，而非空集合 \(P\subset X_m\setminus Y\) 满足：对每个 \(y\in P\) 存在 \(x_y\in Y\) 使
\[
E_{x_y y}\in\mathcal A,
\]
则
\[
\mathcal B(U_{Y\cup P})\subset\mathcal A,
\]
从而 \(Y\) 不可能仍为 maximal coherent grain。也就是说，在真实 fold-gauge 可见矩阵扇区内，一个 fully occupied thin prism 一旦携带通向 maximal grain 的相干 matrix-unit attachment，就必然被吸收到更大的 grain；它不能作为独立 residual prism 存活。
\end{theorem}
\begin{proof}
首先注意，对任意 \(T\in\mathcal B(U_m)\)，由于 \(P_x^+\) 与 \(P_y^+\) 均为一维正交投影，存在唯一标量 \(\lambda_{xy}(T)\in\CC\) 使
\[
P_x^+TP_y^+=\lambda_{xy}(T)E_{xy}.
\]
因此若 \(P_x^+\mathcal A P_y^+\ne0\)，则某个 \(T\in\mathcal A\) 满足 \(\lambda_{xy}(T)\ne0\)，从而
\[
E_{xy}=\lambda_{xy}(T)^{-1}P_x^+TP_y^+\in\mathcal A,
\]
因为 \(P_x^+,P_y^+\in\mathcal A\)。反向显然：若 \(E_{xy}\in\mathcal A\)，则 \(P_x^+\mathcal A P_y^+\ne0\)。

若 \(x=x_0,x_1,\ldots,x_r=y\) 是 \(\Gamma_{\mathcal A}\) 中的一条路径，则由上一段可知每个 \(E_{x_{j-1}x_j}\in\mathcal A\)，并且
\[
E_{x_0x_1}E_{x_1x_2}\cdots E_{x_{r-1}x_r}=E_{xy}.
\]
故同一连通分支内的全部 matrix units \(E_{xy}\) 都属于 \(\mathcal A\)。于是
\[
\mathcal B(U_C)=\Span\{E_{xy}:x,y\in C\}\subset\mathcal A
\]
对每个分支 \(C\) 成立。

另一方面，若 \(x,y\) 属于不同连通分支，则按图的定义必有 \(P_x^+\mathcal A P_y^+=0\)。因此任意 \(T\in\mathcal A\) 在不同分支之间的矩阵块都为零，故
\[
\mathcal A\subset\bigoplus_{C\in\mathcal C(\mathcal A)}\mathcal B(U_C).
\]
与上一段的反向包含合并，得到所示块分解。

最后设 \(Y\in\mathcal C(\mathcal A)\)，且 \(P\subset X_m\setminus Y\) 满足题设 attachment 条件。对每个 \(y\in P\)，由 \(E_{x_y y}\in\mathcal A\) 可知 \(x_y\) 与 \(y\) 在 \(\Gamma_{\mathcal A}\) 中相邻；而 \(x_y\in Y\)，故 \(y\) 与整个分支 \(Y\) 位于同一个连通分支。于是 \(Y\cup P\) 包含在某个连通分支 \(C'\) 中。由已证的分解，\(\mathcal B(U_{C'})\subset\mathcal A\)，从而特别有 \(\mathcal B(U_{Y\cup P})\subset\mathcal A\)。若 \(P\ne\varnothing\)，则 \(C'\) 严格大于 \(Y\)，这与 \(Y\) 作为 maximal coherent grain 相矛盾。

该论证只使用 fold-aware 的纤维平均向量 \(u_x\) 与真实可见矩阵扇区 \(U_m\)。当某些 \(d_x=1\) 时，\(P_x^-=0\) 仅意味着相应暗标准分量为空；上述 rank-one matrix-unit 论证不受影响。证毕。
\end{proof}
% --- End distillation writeback ---

